\documentclass{situationist-preprint}

\title{Failed Intentions and Interpretation}
\msid{AES-09-b}
\msfamily{AES-09}
\msdiscipline{Aesthetics}
\mstype{Consequence}
\mskeywords{intentionalism; failed intentions; interpretation; biography; critical evaluation}

\begin{document}
\maketitle

\begin{abstract}
AES-09-a argues that an author's intentions fix a work's meaning when and only when they are successful, and that failed intentions belong to biography. It may seem to follow that failed intentions are of no use to the interpreter. We derive the opposite consequence in a precise form. Failed intentions cannot fix meaning, but they remain evidence: about the author, about the history of the work's making, and about the achievement the work embodies. We distinguish three uses that interpretation and criticism may legitimately make of them, and one use they may not. The result clarifies the scope of the success account and answers the charge that it renders the study of authors irrelevant to the study of works.
\end{abstract}

\section{The consequence}

On the account of AES-09-a, a semantic intention is successful when its intended meaning is a permissible reading of the work and the work gives its intended audience sufficient grounds to recognize it as intended (Definition~2.4 of AES-09-a). Theorem~3.3 of AES-09-a states that intentions fix meaning when and only when they are successful, and that failed intentions are facts about the author. Corollary~3.4 adds that external evidence, such as letters and statements, bears on interpretation only as evidence about which permissible reading was successfully intended.

The phrase \enquote{belong to biography} can be misread as \enquote{belong nowhere in criticism}. We show that it has a narrower and more useful content.

\begin{proposition}\label{p:uses}
A failed semantic intention is irrelevant to what a work means, but it is relevant (i) to the explanation of the work's features, (ii) to the assessment of the work as an achievement, and (iii) to the identification of which permissible reading a partly successful intention selects. It is not relevant to (iv) the attribution to the work of any content that is not a permissible reading of it.
\end{proposition}

\begin{proof}[Argument]
(iv) is Lemma~3.1 of AES-09-a together with Principle~P1 of that manuscript: a work cannot mean what is not a permissible reading of it, and a failed intention, as such, adds no permissible reading. (i) Features of a work are products of its making, and the intentions with which it was made, successful or not, are among their causes. A failed attempt at irony explains why a passage reads as earnest but strained; the failed intention is part of the explanation of the strain, though not of its meaning. Such explanations are biographical in the sense of AES-09-a, but criticism has always included them. (ii) The artistic value of a work includes the value of the achievement it embodies (AES-04). What a work achieves is measured partly against what it attempted, and an attempt that failed is relevant to that measure. (iii) Intentions are not all-or-nothing. An author may intend a passage to be ironic and to allude to a particular source; if the irony is recognizable and the allusion is not, the intention succeeds in part. Evidence of the whole intention, including its failed part, helps the interpreter to identify the part that succeeded, which is what Corollary~3.4 of AES-09-a permits.
\end{proof}

\section{Two cases}

An author writes an elegy intending its final stanza to undercut the consolation of the preceding ones. Readers, competent in the conventions of the form, take the stanza as a further consolation, and nothing in its wording or in the tradition marks it otherwise. A letter later surfaces stating the intention. On the success account the stanza does not undercut the consolation, since the undercutting reading is not one the work supports. The letter is nonetheless valuable: it explains a hesitancy in the stanza's diction that had puzzled readers, and it tells us what the author was attempting, which bears on how we assess the elegy as an accomplishment. A critic who reports the letter and declines to reinterpret the stanza is doing exactly what Proposition~\ref{p:uses} prescribes.

In a second case, a novelist intends a character's name to allude to a classical figure and to signal the character's fate. The allusion is recognizable to the novel's intended audience; the signal of fate depends on an obscure variant of the myth that the audience cannot be expected to know. The allusion is part of the novel's meaning; the prophecy is not. Evidence that the novelist intended both helps a reader confirm that the allusion is intended rather than accidental, which is the use (iii) allows.

\section{Objections}

\begin{objection}
If failed intentions bear on artistic value, they bear on criticism of the work, and the success account has readmitted them by the back door.
\end{objection}

\begin{reply}
Criticism of a work includes interpretation and evaluation. The success account is a thesis about interpretation: about what a work means. Its consequence for evaluation is not that failed intentions are excluded, but that they enter as facts about an attempt and not as contents of the work. The elegy is not credited with an undercutting it does not perform; it may be judged a work that attempted more than it achieved.
\end{reply}

\begin{objection}
Anti-intentionalists such as \citet{WimsattBeardsley1946} would say that even these uses confuse the poem with its origins.
\end{objection}

\begin{reply}
The confusion they warned against is the identification of a work's meaning with its author's intended meaning. Proposition~\ref{p:uses}(iv) excludes it. What remains, explanation and the assessment of achievement, is the study of a work as the product of an act, which even a strict anti-intentionalist about meaning has no reason to forbid \citep{Levinson1996}.
\end{reply}

\section{Conclusion}

Failed intentions belong to biography, and biography is not idle. It explains features of works, measures their achievement, and helps to locate the part of an intention that succeeded. What it cannot do is make a work mean what the work cannot support.

\begin{thebibliography}{9}\small
\bibitem[Levinson(1996)]{Levinson1996} Levinson, J. (1996). Intention and interpretation in literature. In \emph{The Pleasures of Aesthetics: Philosophical Essays}. Cornell University Press.
\bibitem[Wimsatt and Beardsley(1946)]{WimsattBeardsley1946} Wimsatt, W.~K. and Beardsley, M.~C. (1946). The intentional fallacy. \emph{The Sewanee Review}, 54(3), 468--488.
\end{thebibliography}

\end{document}
